\subsection{Verification}
\label{sec:interference-verification}

Compliance is demonstrated in three stages, and a specification names all
three.

At the design stage the supplier submits an interference study according to
CIGRE TB~553~[30]: the calculated noise spectrum at the AC and DC line
terminals in the band 30\,kHz to 500\,kHz, the calculated radio interference
field strength at the station fence in the band 150\,kHz to 1\,GHz, the
calculated $I_{eq}$ of the DC line, and the design of the PLC filters and the
valve-hall shielding. A preliminary study belongs to the proposal, since the
PLC filter is part of the price; the final study belongs to the design
review.

At the factory stage the RIV of the switchyard apparatus is verified by the
type tests of the product standards, with the test voltage corrected to the
site elevation as in Section~\ref{sec:ri}.

At site the background level is measured at the defined points before first
energisation and the measurement is repeated with the station energised at
0\,\%, 25\,\%, 50\,\%, 75\,\% and 100\,\% of rated power in both power
directions, in fair weather, according to IEEE Std~430~[38] for the field
strength and with coupling capacitors at the line terminals for the PLC-band
noise. The background measurement is what makes the result attributable: in
an interconnection area the natural noise of the surrounding lines can exceed
a limit on its own, and the criterion then becomes that the station raises
the measured level by not more than 3\,dB. The measurement report is part of
the site acceptance documentation.
\FloatBarrier
